\documentclass[12pt]{article}
\usepackage[bookmarks=false]{hyperref}
\usepackage{enumerate}
\usepackage{amssymb,amsmath,amsthm}
\usepackage{graphicx}
\usepackage{geometry}

\pagestyle{empty}

\begin{document}
	
\section*{Discrete version}
Assume that $k$ targets have been identified at radial positions $\theta_i$ with signal strength $s_i$. Let $\vec{p}_i = (\theta_i,s_i)$. A consensus direction will be chosen from among these targets.

Let $0\leq n_i\leq 1$ be the fraction of spins active toward each $\vec{p}_i$. The consensus direction will be chosen based on these spins:
\begin{equation}
\Theta_{foc} = \sum_{i=1}^{k}\theta_in_i\label{eqn:Theta}
\end{equation}

Under the hood, we will assume a Hamiltonian of the form
\[
H = -\frac{k}{N}\sum_{i\neq j}J_{ij}\sigma_i\sigma_j
\]
where $\sigma_i$ is one of $N$ total spins in the system among $k$ groups (directions). This is exactly the same as the Gorbonos PNAS paper. However, we will choose a slightly different interaction kernel $J_{ij}$.
\[
J_{ij} = s_i\cos\left(\pi\left(\frac{|\theta_{ij}|}{\pi}\right)^\nu\right)
\]
If all signal strengths in all directions are equal, we recover the PNAS model. $\nu=0.5$ was used in the PNAS paper and is similar to Mexican hat (see Fig. S1, Table S1, and Fig. 1 for comparison with experiments).

This change easily trickles through the entire framework. To re-purpose Gorbonos's notation, let $\hat{p}_i$ be a unit vector in the direction of $\theta_i$, and let $\hat{V}_{foc}$ be a unit vector in the direction of $\Theta_{foc}$.
Then
\[
\frac{dn_i}{dt} = \frac{1}{k\left(1+\exp\left(-\frac{2ks_i\hat{V}\cdot\hat{p}_i}{T}\right)\right)} - n_i
\]
where the dot product is shorthand notation for the cosine part of $J_{ij}$ and $T$ is the temperature. Assume a separation of time-scales where neural updates of spins occurs far faster than the physical rotation of an insect. In that case, we can treat the radial positions $\theta_i$ of the targets as roughly constant on the time-scale of updating the mental torque (consensus direction). Then, taking the derivative of Eqn. \eqref{eqn:Theta} and plugging in, we get
\begin{equation}
	\frac{d\Theta_{foc}}{dt} = \sum_{i=1}^{k}\theta_i\frac{dn_i}{dt} = \frac{1}{k}\sum_{i=1}^{k}\frac{\theta_i}{\left(1+\exp\left(-\frac{2ks_i\hat{V}\cdot\hat{p}_i}{T}\right)\right)} - \Theta_{foc}
\end{equation}
This can be interpreted directly as a torque model by interpreting $\Theta_{foc}$ as the current direction the locust is facing. Also, $\theta_i - \Theta_{foc}$ is then the angle of target $i$ relative to the current heading of the locust, $\theta_{i,foc}$ and using an expanded definition of dot product based on interaction kernel $J_{ij}$,
\[
s_i\hat{V}\cdot\hat{p}_i = J_{i,foc}.
\]

\subsection*{Signal strengths}
This discrete version assumes that radial target locations are identified by a separate perception process. If signal strength does not vary with distance, then the underlying assumption is that all targets are delta functions. Avoiding this, the best choice for signal strength is probably some measure of the perceived angular width of the target. The problem of not being able to tell two overlapping locusts from one close locust is avoided by the prior target identification, and size scales with distance in the proper manner (closer yields larger numbers while staying bounded) and is also biological in the sense that visual size of the object IS the actual signal reaching the eye.

\section*{Continuous version}
An issue here is that since there are angles that don't have targets, you need this $\eta$ term in Gorbonos [41] to inhibit angles which don't have signal. Given that we are supposed to take $\eta N >> J$ and $N\rightarrow\infty$, maybe it isn't important, but it's hard to say.

\end{document}